\scanpage{52}

À \deleted{\unsure{symét}} op. près de $\unsure{\mathfrak S_3|\mathcal S_3}$, \unsure{SPS|PS|SS|OPS} donc que le système est de la forme $(\rho,\rho',\rho'')$, avec $\rho,\rho'$ d'ordre $3$, donc chacun de la forme $\rho_s$ ou $\rho_f$

\unsure{1)|A)} $\rho=\rho_f,\ \rho'=\rho_s$ il f. $s\prec f$ (car si $s$ opposé à $f$, $\rho_s=\rho_f^{-1}$ donc $\rho,\rho'$ n'engendrent pas $H$, donc on est dans le cas canonique

\margin{$12$}

\unsure{2)|B)} $\rho=\underset{\substack{\shortparallel\\\rho'_{f'}}}{\rho_s},\quad\rho'=\underset{\substack{\shortparallel\\\rho'_{s'}}}{\rho_f}$ il f. $s\unsure{\prec|\nprec} f$ i.e. $s'\prec f'$, donc on est dans le cas anticanonique

\margin{$12$}

3) $\rho=\rho_s,\ \rho'=\rho_{\deleted{\unsure{s'}}t}$ \inserted{$s,t\in S$} $(s\ne t)$ on est dans le \inserted{$1\textsuperscript{er}$} cas de e)

\margin{$12$}

4) $\rho=\underset{\substack{\shortparallel\\\rho'_{s'}}}{\rho_f},\quad\rho'=\underset{\substack{\shortparallel\\\rho'_{t'}}}{\rho_{\deleted{f'}g}}$ \inserted{$f,g\in F$} $(f\ne\deleted{f'}g)$ on est dans le $1\textsuperscript{er}$ cas de e)

\margin{$12$}

\uline{Corollaire} Les types numériques possibles des syst. de Schwartz pour $H\simeq\unsure{A_4|\mathcal A_4|\mathfrak A_4}$ sont

\[
\begin{aligned}
\text{a)}\ &(3,3,2)\ (3,2,3)\ (2,3,3)\qquad\text{[cas canonique et anticanonique]}\\
\text{b)}\ &(3,3,3)
\end{aligned}
\]

Et la classe de conjugaison sous $\widetilde H=\mathrm{Aut}_{\mathrm{gr}}(H)$. $(\simeq\mathfrak S_4)$ est déterminée par le type numérique.

Dans un $H$ $(\simeq A_4)$ un syst. de Schwartz est d'un des types suivants (où $H=\mathrm{Aut}^{+}(\Pi)$, $\Pi$ un tétraèdre comb., dont on n'a pas d'orientation privilégiée)

\margin{\uline{NB} La donnée de $H$ équivaut à celle de $\Pi$ cf plus haut}

a) On a un repère $r$ (\unsure{uniquemt|uniqmt|unique} déterminé) de $\Pi$, et $\Sigma=\Sigma_r$ (cas canonique)

\deleted{b) On a un \deleted{\unsure{orientation}} repère $r=\{s,\{s,s'\},\{s,s',s''\}\}$ bien déterminé et $\Sigma=(\rho_s^r,\rho_{s'}^r,\rho_{s''}^r)$, où l'exposant}

% archive sheet number: 51
